\begin{lemma}\label{minremtab}
Every nonempty skew shape \(\tau\) has a minimal \(\ssee\)-removable ribbon tableau \((\Lambda,\ttt)\).
\end{lemma}
\begin{proof}
We go by induction on \(|\tau|\). If \(|\tau| = 1\), then we trivially take \(\Lambda = \{\tau\}\) and \(\ttt(1) = \tau\). Now let \(|\tau|>1\) and make the induction assumption on skew shapes \(\tau'\) with \(|\tau'| <|\tau|\).
By Lemma~\ref{minremexist}, \(\tau\) has a minimal \(\ssee\)-removable ribbon \(\xi\) such that \(\cont(\xi) \in \Psi\). Let \(\tau' = \tau \backslash \xi\). If \(\tau' = \varnothing\), we are done, taking \(\Lambda = \{\xi\}\), \(\ttt(1) = \xi\). Assume \(\tau' \neq \varnothing\). Then \(\tau'\) is a skew shape by Lemma~\ref{tabisskew}, which by the induction assumption has a minimal \(\ssee\)-removable ribbon tableau \((\Lambda', \ttt')\). But then taking \(\Lambda = \Lambda' \cup \{\xi\}\), \(\ttt(i) = \ttt'(i)\), \(\ttt(|\Lambda|) = \xi\), for all \(i \in [1,|\Lambda'|]\), we have that \((\Lambda, \ttt)\) is a minimal \(\ssee\)-removable ribbon tableau for \(\tau\).
\end{proof}

\begin{lemma}\label{minremtabisKos}
If \((\Lambda,\ttt)\) is a minimal \(\ssee\)-removable ribbon tableau for a skew shape \(\tau\), then \((\Lambda,\ttt)\) is a cuspidal Kostant tableau for \(\tau\).
\end{lemma}
\begin{proof}
By Lemma~\ref{minremiscusp}, every \(\lambda \in \Lambda\) is cuspidal, so it remains to verify that \((\cont(\ttt(i)))_{i=1}^{|\Lambda|}\) is a Kostant sequence by showing that \(\cont(\ttt(j)) \succeq \cont(\ttt(j+1))\) for all \(j \in [1,|\Lambda|-1]\). Indeed, \(\ttt(j+1)\) is a minimal \(\ssee\)-removable ribbon tableau in \(\bigsqcup_{i=1}^{j+1}\ttt(i)\), and \(\ttt(i)\) is a minimal \(\ssee\)-removable ribbon tableau in \(\bigsqcup_{i=1}^{j} \ttt(i)= (\bigsqcup_{i=1}^{j+1}\ttt(i)) \backslash \ttt(j+1)\), so the result follows by Lemma~\ref{remsinorder}.
\end{proof}

\begin{lemma}\label{cuspisribbon}
Every cuspidal skew shape is a ribbon.
\end{lemma}
\begin{proof}
Let \(\tau\) be a cuspidal skew shape, so that \(\cont(\tau) \in \Phi_+\), and assume by way of contradiction that \(\tau\) is not a ribbon. By Lemma~\ref{minremtab}, there exists a minimal \(\ssee\)-removable ribbon tableau for \(\tau\). As \(\tau\) is not itself a ribbon, it must be that \(|\Lambda|>1\). By Lemma~\ref{minremtabisKos} we have 
\begin{align*}
\cont(\ttt(1)) \succeq \cdots \succeq \cont(\ttt(|\Lambda|)).
\end{align*}
As
\(
\cont(\ttt(1)) + \cdots + \cont(\ttt(|\Lambda|)) = \cont(\tau), 
\)
we have by Lemma~\ref{GenConvexLem} that
\begin{align*}
\cont(\ttt(1)) \succeq \cont(\tau)\succeq \cont(\ttt(|\Lambda|)).
\end{align*}
But \(\ttt(|\Lambda|)\) is a minimal \(\ssee\)-removable ribbon in \(\tau\), so \(\cont(\ttt(|\Lambda|)) \succ \cont(\tau)\) by Lemma~\ref{shapehooktest}, giving the desired contradiction.
\end{proof}

\begin{lemma}\label{cuspisindiv}
If \(\xi\) is a cuspidal skew shape, then \(\cont(\xi) \in \Psi\).
\end{lemma}
\begin{proof}
By Lemma~\ref{cuspisribbon} we have that \(\xi\) is a ribbon. If \(\cont(\xi) \in \Phi_+ \backslash \Psi\), then \(\cont(\xi) = m\delta\) for some \(m >1\). But then \(\xi\) has an \(\ssee\)-removable skew hook \(\nu \subsetneq \xi\) of content \(\cont(\xi) \succeq \delta\), a contradiction of Lemma~\ref{CuspForSkewHook}, so \(\cont(\xi) \in \Psi\).
\end{proof}
